% !PRINTED 647
\weberpage{647}
%<<CONT
Es ist aber $r^s \pm r^{s'} = r^s (1 \pm r^{h-2s}) = r^s (1 \mp r^{\mu+h-2s})$ immer durch $1 - r$ theilbar, und daraus würde folgen, dass $\mathfrak{E}(r)$ durch $1 - r$, was keine Einheit ist, theilbar wäre, was dem Begriffe der Einheit widerspricht. Also ist die Annahme unzulässig, dass in der Formel (4) der Exponent $h$ ungerade sei, und wir können wie oben
\[
\mathfrak{E}(r) = r^{\tfrac12 h}\, e(r)
\]
setzen, worin $e(r)$ eine reelle Einheit ist. Damit ist der Satz 1. bewiesen.

Wir wollen ein System von reellen Einheiten hervorheben: Nach §. 169 sind die $\mu$ Grössen $1 - r^n$ alle unter einander associirt. Demnach ist der Quotient
\[
\frac{1 - r^n}{1 - r^{n'}},
\]
worin $n, n'$ zwei relative Primzahlen zu $m$ bedeuten, eine Einheit. Daraus leitet man aber, wenn
\[
r = e^{\frac{2\pi i}{m}}
\]
gesetzt wird, die reelle Einheit
\weq{6}{\dfrac{\sin \dfrac{n\pi}{m}}{\sin \dfrac{n'\pi}{m}} = \dfrac{r^{\frac{n}{2}} - r^{-\frac{n}{2}}}{r^{\frac{n'}{2}} - r^{-\frac{n'}{2}}} = r^{\frac{n'-n}{2}}\, \dfrac{1 - r^n}{1 - r^{n'}}}
her. Ist $m$ gerade, so kann man $n' = n + \dfrac{m}{2}$ setzen, und erhält für diesen Fall die reellen Einheiten
\weq{7}{\tau_n = \mathrm{tang}\,\dfrac{n\pi}{m},}
worin $n$ jede ungerade Zahl bedeuten kann.

Ersetzt man $n$ durch $\tfrac{1}{2} m - n$, so geht $\tau_n$ in den reciproken Werth über. Lässt man $n$ die Reihe der Zahlen
\[
1,\ 3,\ 5,\ \cdots,\ \frac{m}{4} - 1
\]
durchlaufen, so erhält $\tau_n$ lauter positive echt gebrochene Werthe.

\begin{center}
\rule{0.2\linewidth}{0.4pt}
\end{center}
